% ShelfSense weekly report 5 (DRAFT by the coding assistant from the week's runs; the team checks and edits it).
% Paste into the Overleaf main file as the next "\section{Weekly Report 5}". Upload figures/r5_*.png to the
% Overleaf project (they are not in Git: the screenshots show the NutriCoco demo). Needs \usepackage{xcolor} and
% \newcommand{\TODO}[1]{\textcolor{red}{[#1]}} in the preamble. Lines marked \TODO{...} need a team member.

\section{Weekly Report 5}

\subsection{Follow-up from last week}
Status of the targets set in Weekly Report 4:
\begin{enumerate}
  \item \textbf{Correct the pre-labels of the 222 keyframes in Roboflow --- \TODO{status}.}
  \item \textbf{Fine-tune the detector on our frames and report mAP and counting error --- \TODO{status}.}
  The training script is ready; it needs the corrected labels of item~1.
  \item \textbf{Re-score the baseline with the confirmed Mojo label --- \TODO{status}.}
  \item \textbf{Record shop fridges with NutriCoco's paid front zone --- Changed.} NutriCoco will send
  photos and videos of their cups in shop fridges themselves; we wrote a photo guide (English and
  Bangla) so that the material is useful for training (Section~\ref{sec:r5-ai}).
  \item \textbf{First KPI checks with evidence --- Partly done.} The paid-placement check against each
  shop's agreement is implemented and tested (Section~\ref{sec:r5-ai}); it switches on as soon as the
  detector knows NutriCoco's cups. Brand reading on crops is not yet in the audit pipeline.
  \item \textbf{Re-use check robust to cropping --- Done.} All 10\%-cropped copies are now caught and
  no honest re-recording of the same shelf is flagged (Section~\ref{sec:r5-reuse}).
\end{enumerate}

\subsection{Work narrative of the current week}
Our industry partner NutriCoco met the team and asked for much more than a fridge check: tracking from
the factory to the shelf, delivery staff performance, stock and restocking, Pandamart purchase orders,
reseller ordering and live location sharing. This week we turned the audit pipeline into a working
system with a phone app for the delivery staff, office screens for NutriCoco and a server on our PC.
Everything below was tested on the team's computer with sample shops; 97 automated tests pass.

\subsubsection{Partner requirements and scope}
NutriCoco's nine requirements and their answers to our questions (the same staff deliver and check
fridges; Android and iPhone; agreements per shop must be changeable; motorcycle delivery; most sales go
to Pandamart, which sends purchase orders; dispatch records are on paper; no deadline) are written up in
a requirements summary for NutriCoco to sign. Table~\ref{tab:r5-scope} shows how each requirement is
covered. Only requirement~3 (AI shelf placement) is part of the course topic; the others make the
system usable for the partner. Phase~1 (this semester) covers the delivery and audit flow and the
NutriCoco cup detector; phase~2 covers an internet server, native apps and shop-owner ordering.

\begin{table}[!t]
\centering
\caption{NutriCoco's requirements and their status (demo on the team's PC, sample data).}
\label{tab:r5-scope}
\begin{tabular}{p{0.30\linewidth}p{0.40\linewidth}p{0.20\linewidth}}
\hline
Requirement & What the system does & Status \\
\hline
1 Factory-to-shelf tracking & production $\rightarrow$ motorcycle $\rightarrow$ shop delivery $\rightarrow$ fridge video & working \\
2 Stock and restocking & factory and motorcycle stock, restock alerts & working \\
3 AI placement and KPIs & cups vs.\ agreement, placement, gaps, before/after & waits for photos \\
4 Pandamart & POs entered by hand, over/under delivery, returns & no API access \\
5 Staff performance & check-in/out, time per shop, km, fuel cost & working \\
6 Live location & live map, expiring share links, arrival alerts & app must be open \\
7 Wholesale vs.\ reseller & customer types, separate sections and prices & working \\
8 Direct ordering & reseller ordering page with PIN & working \\
\hline
\end{tabular}
\end{table}

\subsubsection{Server and phone app}
A small server on our PC (Python's built-in HTTP server, nothing extra to install) receives videos,
runs the Sprint~5 audit pipeline in a background worker and stores everything in SQLite. Phones need
HTTPS for camera and GPS, so the demo is reached through a free Cloudflare tunnel. Each rider and the
office log in with a 6-digit PIN; only salted PBKDF2 hashes are stored and five wrong PINs lock the
address for five minutes. Riders see only their own data.

The rider app is an installable web app (one code base for Android and iPhone, English and Bangla,
Fig.~\ref{fig:r5-app}). It records the fridge video \emph{inside} the app at 720p: a live recording
cannot be an old video from the gallery, and an 8\,s test recording was 414\,kB instead of the 83\,MB
of a 4K 60\,fps phone video sent by a team member. Videos, GPS points and deliveries made without a
signal wait in an offline outbox and are sent later with their original time.
Table~\ref{tab:r5-uploads} shows real uploads through the server.

\begin{table}[!t]
\centering
\caption{Videos uploaded through the server and audited (CPU detector).}
\label{tab:r5-uploads}
\begin{tabular}{lccl}
\hline
Video & Size & Time (s) & Result \\
\hline
WA0089 (supermarket, 576$\times$1024) & 3.3\,MB & 8.5 & re-audit: same footage as an earlier audit \\
WA0061 (supermarket, 576$\times$1024) & 0.8\,MB & 3.6 & re-audit: same footage as an earlier audit \\
Team phone video (3840$\times$2160, 60\,fps) & 83\,MB & 16.8 & review: no products in 6 of 10 keyframes \\
In-app recording (test camera, 1280$\times$720) & 0.4\,MB & --- & re-audit: filmed 1.6\,km from the chosen shop \\
\hline
\end{tabular}
\end{table}

Both re-audit results on supermarket videos are correct: the same files had been audited before, and
the re-use check found them. The phone video showed a person in a room, not a shelf; the detector found
0--1 ``products'' per frame (once a face) and the audit asked for a review. Every video is now also
checked for \emph{where} and \emph{when} it was filmed: more than 200\,m from the shop gives a warning,
more than 1\,km a re-audit, and a video recorded more than 6\,h before it was sent a warning.

\subsubsection{Re-used footage that is cropped or resized}\label{sec:r5-reuse}
The keyframe pHash of Report~4 caught re-encoded and half-size copies but missed a 10\% crop. A second
check now matches ORB keypoints between keyframes, fits a homography with RANSAC and warps the old
keyframe onto the new one. Matching points alone cannot decide, because riders honestly film the same
shelf again; the decision therefore uses the pixel correlation after alignment, which is almost 1 only
for a copy. We tested it on altered copies of three test videos (positives) and on all 24 pairs of
neighbouring test clips, which were filmed minutes apart and often show the same shelf (negatives).
Copies correlated 0.994--1.000 after alignment; honest re-recordings reached at most 0.964 on single
keyframes, so the threshold is 0.98 (Table~\ref{tab:r5-reuse}). To keep the audit fast, the check runs
only when the pHash finds nothing and only on the five earlier audits with the closest pHashes (the
original ranked first or second of 124 for every crop): an audit with a cropped copy takes 4.5\,s, an
honest one 2.3--2.5\,s as before.

\begin{table}[!t]
\centering
\caption{Re-use check on altered copies and honest neighbouring clips (keyframes or videos flagged).}
\label{tab:r5-reuse}
\begin{tabular}{lcc}
\hline
 & pHash only (Report 4) & pHash + local features \\
\hline
Re-encoded copies (3 videos) & 3 / 3 & 3 / 3 \\
Half-size copies (3 videos) & 3 / 3 & 3 / 3 \\
10\% cropped copies (3 videos) & 0 / 3 & 3 / 3 \\
Copied keyframes recognised & --- & 88 / 90 \\
Honest neighbouring clips flagged (24 pairs) & 0 / 24 & 0 / 24 \\
Honest keyframes flagged & --- & 0 / 255 \\
\hline
\end{tabular}
\end{table}

\subsubsection{Delivery tracking}
While the app is open, the phone sends a GPS point every 15\,s or 30\,m. A rider is checked in when
inside a shop's 80\,m circle and checked out only beyond 120\,m, so standing at the edge does not
flicker. Points with an accuracy worse than 50\,m or implying more than 90\,km/h are stored but never
change a visit or add distance. A simulated two-hour route through the sample Mirpur shops (167 points)
gave 5.65\,km and 8 visits; one of them was the rider passing a shop without stopping, so stays under
2 minutes now count as ``passed by''. Route reports give start and end time, km, fuel cost (with
placeholder fuel numbers until NutriCoco sets theirs), shops visited and time per shop.

\subsubsection{Distribution records}
Production, loading of each motorcycle, deliveries with returned quantities and reasons, Pandamart
purchase orders and reseller orders are recorded in the app. The office sees over- and
under-delivery per PO (the demo PO was 6 cups short), stock in the factory and on each motorcycle,
returns by reason and restock alerts. Resellers order on their own page at wholesale prices.

\subsubsection{AI parts}\label{sec:r5-ai}
\textbf{Agreement check.} For each fridge video the system takes the agreement valid on the day it was
recorded (cups, shelf, position) and reports the cups visible in the best keyframe, the cups in the
agreed place, empty gaps on the agreed shelf, the share of that shelf, and the change since the shop's
previous video. Following our precision-over-recall rule, a shortfall is reported as a violation only
when the video passed the coverage checks; otherwise it is ``unclear'' and a person decides. The check
is tested on hand-made detections and stays ``pending'' until the detector has a NutriCoco cup class.
The training script for that detector splits photos by shop visit (near-identical photos of one fridge
must not be in both training and test) and reports mAP per class and the cup-counting error; it was
tested end to end on a small artificial dataset.

\textbf{Audit assistant.} The ``Ask about this audit'' chat runs a free open model on our PC through
Ollama, so no data leaves the PC. It only receives one audit report and strict rules (use only this
data, cite the keyframe, say when the data does not show something, never decide payments, answer in
the language of the question). Table~\ref{tab:r5-assistant} compares three models on 10 questions
(6 English, 4 Bangla) about 3 real audits. With the first rules the models answered Bangla questions in
English and gave payment advice; repeating the language rule next to the question and fixing the
closing sentence for payment questions raised Gemma~3 4B from 25 to 29 of 30. We chose Gemma~3 4B.

\begin{table}[!t]
\centering
\caption{Audit assistant: checks passed out of 30 answers per model (rules v2; v1 in brackets).}
\label{tab:r5-assistant}
\begin{tabular}{lccc}
\hline
 & Gemma 3 4B & Qwen2.5 3B & Qwen2.5 1.5B \\
\hline
All checks & 29 (25) & 22 (24) & 22 (19) \\
Facts correct (of 15) & 14 (15) & 10 (11) & 11 (11) \\
Refuses what the data does not show (of 9) & 9 (9) & 6 (9) & 5 (8) \\
Leaves payment to the auditor (of 6) & 6 (1) & 6 (4) & 6 (0) \\
Bangla questions answered in Bangla (of 12) & 9 (1) & 12 (4) & 12 (9) \\
Median time per answer (s, GTX 1650) & 11.3 & 3.4 & 2.9 \\
\hline
\end{tabular}
\end{table}

\begin{figure}[!t]
\centering
\includegraphics[width=0.32\linewidth]{figures/r5_field_app.png}\hfill
\includegraphics[width=0.64\linewidth]{figures/r5_office_live.png}
\caption{Left: rider app in Bangla (tracking, in-app fridge video, delivery). Right: NutriCoco's live
map with shops, their check-in circles and the check-in feed (sample data).}
\label{fig:r5-app}
\end{figure}

\subsubsection{Challenges}
\begin{itemize}
  \item \textbf{Background tracking.} A web app can only track while it is open on the screen; the app
  keeps the screen on while tracking. Tracking with a locked phone needs a native app (phase~2).
  \item \textbf{Bugs found by testing.} A single wrong GPS point once checked a rider out of a shop;
  report days were counted in UTC instead of Bangladesh time; a video waited forever for a GPS position
  when the permission question was not answered; the app did not know about a visit that started
  before it was opened. All four are fixed and covered by tests.
  \item \textbf{No NutriCoco photos yet.} Cup counting depends on them; the rest of the system was
  built so that the trained detector only has to be dropped in.
  \item \textbf{Pandamart.} No public API is available, so purchase orders are entered by hand.
\end{itemize}

\subsection{Target for next weekly update}
\begin{enumerate}
  \item Get the requirements summary signed and the pilot shop list with each shop's agreement from
  NutriCoco; replace the sample shops and prices.
  \item Label NutriCoco's photos (two members check every test photo), train the cup detector and
  report mAP and the cup-counting error on unseen photos.
  \item Pilot ride in Mirpur~10/11 with a team member's phone: measure check-in accuracy, false visits
  and the km against the motorcycle's odometer.
  \item Carried over: correct the 222 keyframes in Roboflow and fine-tune the shelf detector.
  \item Test the re-use check on a real case: the same fridge filmed on two different days by a team
  member, and an old video of it cropped on a phone.
\end{enumerate}

\subsection{Individual contribution}
\begin{table}[!h]
\centering
\caption{Contribution summary.}
\begin{tabular}{lcc}
\hline
Task List & Member Assigned & Comments \\
\hline
NutriCoco meeting, requirements and NDA & \TODO{member} & 9 requirements \\
Audit assistant evaluation & \TODO{member} & 3 models, 30 answers each \\
Server, login and tracking & \TODO{member} & geofence check-in/out \\
Rider app and offline outbox & \TODO{member} & English / Bangla \\
Office screens and reseller page & \TODO{member} & orders, stock, POs \\
Agreement check and training script & \TODO{member} & waits for photos \\
Re-use check for cropped copies & \TODO{member} & 9/9 copies, 0/24 false \\
Testing on phones & \TODO{member} & 4K video, live location \\
\hline
\end{tabular}
\end{table}
